\documentclass[12pt,a4paper]{report}

% ── Encoding & Fonts ──
\usepackage[utf8]{inputenc}
\usepackage[T1]{fontenc}
\usepackage{lmodern}

% ── Page Layout ──
\usepackage[margin=1in,headheight=15pt]{geometry}
\usepackage{setspace}
\onehalfspacing

% ── Language ──
\usepackage[english]{babel}

% ── Graphics & Colors ──
\usepackage{graphicx}
\usepackage[dvipsnames,svgnames,table]{xcolor}

% ── Math ──
\usepackage{amsmath,amssymb,amsfonts}

% ── Tables ──
\usepackage{booktabs}
\usepackage{tabularx}
\usepackage{longtable}
\usepackage{multirow}

% ── Lists ──
\usepackage{enumitem}

% ── Code Listings ──
\usepackage{listings}
\lstset{
  basicstyle=\ttfamily\small,
  keywordstyle=\color{blue}\bfseries,
  commentstyle=\color{gray},
  stringstyle=\color{red},
  breaklines=true,
  frame=single,
  backgroundcolor=\color{gray!5},
  numbers=left,
  numberstyle=\tiny\color{gray},
  numbersep=8pt,
  tabsize=4,
  showstringspaces=false,
}

% ── Hyperlinks ──
\usepackage[colorlinks=true,linkcolor=NavyBlue,citecolor=Green,urlcolor=RoyalBlue]{hyperref}

% ── Headers & Footers ──
\usepackage{fancyhdr}
\pagestyle{fancy}
\fancyhf{}
\fancyhead[L]{\leftmark}
\fancyhead[R]{\thepage}
\renewcommand{\headrulewidth}{0.4pt}

% ── Chapter Style ──
\usepackage{titlesec}
\titleformat{\chapter}[display]
  {\normalfont\huge\bfseries}
  {\chaptertitlename\ \thechapter}{20pt}{\Huge}
\titlespacing*{\chapter}{0pt}{-20pt}{30pt}

% ── Caption Styling ──
\usepackage{caption}
\captionsetup{font=small,labelfont=bf}

% ── Misc ──
\usepackage{parskip}
\usepackage{float}

% ════════════════════════════════════════════════
%  DOCUMENT METADATA
% ════════════════════════════════════════════════
\title{\textbf{OpenFOAM Simulation Guide}\\[6pt]
       \Large Hydrogen--Methane T-Junction Mixing\\
       Bottom-Injection Parametric Study (RANS $k$-$\omega$ SST)}
\author{Yash}
\date{\today}

\begin{document}

% ── Title Page ──
\maketitle
\thispagestyle{empty}

% ── Table of Contents ──
\tableofcontents
\newpage


% ════════════════════════════════════════════════════════════════
%  CHAPTER 1 — Introduction
% ════════════════════════════════════════════════════════════════
\chapter{Introduction}
\label{ch:intro}

This guide documents the OpenFOAM setup for a \emph{parametric study}
of hydrogen injection into a high-pressure methane pipeline through a
T-junction, with the branch pipe entering from the \textbf{bottom} of
the main pipe ($y = -R_1$).  The setup uses a steady RANS approach
($k$-$\omega$ SST) run as a transient simulation with adaptive
time-stepping to capture the quasi-steady mixing field.

The configuration is derived from the benchmark study of
\textbf{Eames, Austin \& Wojcik} (\emph{International Journal of
Hydrogen Energy}, 47:25745--25754, 2022), adapted for:
\begin{itemize}
  \item \textbf{Bottom injection} --- branch enters at $y = -R_1$
        instead of $y = +R_1$ (top), so that the lighter H$_2$ must
        rise against gravity through the heavier CH$_4$.
  \item \textbf{Parametric design of experiments (DoE)} --- the branch
        diameter ratio $d/D$ and velocity ratio VR are varied
        systematically using a Sliced Latin Hypercube Sampling design.
  \item \textbf{Half-domain symmetry} --- an $x = 0$ symmetry plane
        halves the cell count and computational cost.
  \item \textbf{Shortened domain} --- the upstream main pipe is
        reduced from $10\,D$ to $5\,D$ (total $L = 6.9$\,m) to
        further reduce cost without affecting downstream mixing
        metrics.
\end{itemize}

\section{Physical Problem}

Pure hydrogen is injected at 6.9\,MPa into a methane natural-gas
pipeline through a T-junction at various angles.  Two aspects make
the flow physically rich:

\begin{enumerate}
  \item The density ratio $\rho_{\mathrm{CH_4}}/\rho_{\mathrm{H_2}}
        \approx 8$ at 6.9\,MPa generates strong
        \textbf{buoyancy-driven stratification}.
  \item For bottom injection, buoyancy \emph{assists} vertical
        transport of H$_2$ (lighter gas rises from the bottom),
        whereas for top injection buoyancy \emph{traps} H$_2$ against
        the upper wall.
\end{enumerate}

\section{Objectives}

The parametric study quantifies, for each $(d/D,\,\mathrm{VR},\,\alpha)$
combination:

\begin{itemize}
  \item The \textbf{Coefficient of Variation (CoV)} of $Y_{\mathrm{H_2}}$
        at the outlet --- the primary mixing metric.
  \item The \textbf{pressure drop} $\Delta p$ across the junction ---
        the energy cost of injection.
  \item The \textbf{mass balance error} --- numerical credibility check.
  \item The \textbf{H$_2$ plume centroid migration} --- verification that
        buoyancy acts correctly (plume should rise for bottom injection).
\end{itemize}

\section{Reference Setup}

\begin{table}[H]
\centering
\caption{Comparison of the reference Eames setup and this parametric study.}
\label{tab:comparison}
\small
\begin{tabularx}{\textwidth}{@{}lXX@{}}
\toprule
\textbf{Setting} & \textbf{Eames 2022} & \textbf{This study} \\
\midrule
Solver & \texttt{rhoReactingBuoyantFoam} & \texttt{rhoReactingBuoyantFoam} \\
Turbulence & SA-DDES & $k$-$\omega$ SST (RANS) \\
Injection & Top ($y = +R_1$), $90^{\circ}$ & Bottom ($y = -R_1$), $30$--$150^{\circ}$ \\
Domain & Full ($20\,D$ length) & Half (symmetry at $x = 0$, $15\,D$ length) \\
$d/D$ & Fixed (0.25) & Variable (0.20--0.45) \\
VR & Fixed & Variable (0.7--5.0) \\
Cell count & 2.8--4.8\,M & 0.4--0.5\,M (half-domain RANS) \\
End time & 5\,s & 1.2\,s ($\sim 2$ flow-throughs) \\
\bottomrule
\end{tabularx}
\end{table}


\section{Case Directory Layout}

\begin{lstlisting}[language={},numbers=none,frame=none,
  backgroundcolor=\color{white}]
openfoam_case_rans_doe_<angle>_bottom_base/
  0/                  Initial & boundary conditions
    U  p  p_rgh  T  H2  CH4  k  omega  nut  alphat
  constant/
    g                              gravity vector
    momentumTransport              turbulence (k-omega SST)
    turbulenceProperties           Sct, kMin, omegaMin
    thermophysicalProperties       mixture/species (JANAF)
    combustionProperties           combustion OFF
    reactions                      empty
  system/
    controlDict                    run control + function objects
    fvSchemes                      discretisation (upwind)
    fvSolution                     linear solvers, PIMPLE
    blockMeshDict                  background mesh (half-domain)
    snappyHexMeshDict              conformal mesh
    surfaceFeatureExtractDict      feature extraction
    decomposeParDict               parallel decomposition (scotch)
  generateSTL.py                   parametric STL generator
  doe/
    lhs_design.py                  sliced LHS design generator
    stamp_cases.py                 case materialiser (tokens -> values)
    run_doe.sh                     sequential campaign orchestrator
    doe_design.csv                 design matrix
  tools/
    all_metrics.py                 retro-active metric extraction
    make_figures.py                PyVista figure generation
    make_distance_figures.py       cross-section strip plots
\end{lstlisting}


% ════════════════════════════════════════════════════════════════
%  CHAPTER 2 — Governing Equations
% ════════════════════════════════════════════════════════════════
\chapter{Governing Equations}
\label{ch:goveq}

\section{Continuity}

\[
  \frac{\partial \rho}{\partial t}
  + \nabla \!\cdot\! (\rho \mathbf{u}) = 0
\]

\section{Momentum (Buoyant Pressure Decomposition)}

The solver writes $p = p_{\mathrm{rgh}} + \rho\,\mathbf{g}\!\cdot\!
\mathbf{h}$ and advances $p_{\mathrm{rgh}}$:

\[
  \frac{\partial (\rho u_i)}{\partial t}
  + \frac{\partial (\rho u_i u_j)}{\partial x_j}
  = -\frac{\partial p_{\mathrm{rgh}}}{\partial x_i}
    - (\mathbf{g}\!\cdot\!\mathbf{h}) \frac{\partial \rho}{\partial x_i}
    + \frac{\partial}{\partial x_j}\!
      \left[\mu_e\!\left(\frac{\partial u_i}{\partial x_j}
      + \frac{\partial u_j}{\partial x_i}\right)\right]
    + \rho g_i
\]

This decomposition is essential: removing the hydrostatic component
from the momentum equation prevents spurious numerical noise in the
pressure field for flows where the $\rho_{\mathrm{CH_4}} / \rho_{\mathrm{H_2}}
\approx 8$ density contrast drives buoyancy-dominated stratification.

\section{Species Transport}

For each species $k \in \{\mathrm{H_2},\mathrm{CH_4}\}$:

\[
  \frac{\partial (\rho Y_k)}{\partial t}
  + \nabla\!\cdot\!(\rho Y_k \mathbf{u})
  = \nabla\!\cdot\!\left(\frac{\mu_e}{Sc_t}\, \nabla Y_k\right)
\]

\noindent CH$_4$ is the \textbf{inert (balance) species}:
$Y_{\mathrm{CH_4}} = 1 - Y_{\mathrm{H_2}}$.  The turbulent Schmidt
number $Sc_t = 0.5$ (following Eames et al.).

\section{Sensible Enthalpy}

\[
  \frac{\partial (\rho h)}{\partial t}
  + \nabla\!\cdot\!(\rho \mathbf{u} h)
  + \frac{\partial (\rho K)}{\partial t}
  + \nabla\!\cdot\!(\rho \mathbf{u} K)
  - \frac{\partial p}{\partial t}
  = \nabla\!\cdot\!(\mu_e \nabla h)
\]

with $K = \tfrac{1}{2}|\mathbf{u}|^2$.

\section{Closure: $k$-$\omega$ SST}

The $k$-$\omega$ SST model (Menter 1994) blends $k$-$\omega$ near
walls with $k$-$\varepsilon$ in the freestream.  It provides robust
and accurate predictions of separated and recirculating flows at the
T-junction without requiring the fine temporal resolution of DES/LES.


% ════════════════════════════════════════════════════════════════
%  CHAPTER 3 — Geometry and Meshing
% ════════════════════════════════════════════════════════════════
\chapter{Geometry and Meshing}
\label{ch:mesh}

\section{Geometry Definition}

\begin{table}[H]
\centering
\caption{Geometry dimensions.  The branch diameter $d_2$ and
         injection angle $\alpha$ are DoE variables.}
\label{tab:geom}
\begin{tabular}{@{}llr@{}}
\toprule
\textbf{Component} & \textbf{Parameter} & \textbf{Value} \\
\midrule
Main pipe   & Diameter $D_1$       & 460\,mm \\
            & Length $L$           & 6900\,mm ($= 15\,D_1$) \\
            & Upstream of junction & 2300\,mm ($= 5\,D_1$) \\
            & Downstream           & 4600\,mm ($= 10\,D_1$) \\
            & Axis                 & $Z$-direction \\
Branch pipe & Diameter $D_2$       & $0.20$--$0.45 \times D_1$ (DoE) \\
            & Length               & $\max(1.38,\, R_1 + 12\,D_2)$\,m \\
            & Junction centre      & $Z_{\mathrm{jct}} = 2.3$\,m \\
            & Injection angle      & $\alpha = 30^{\circ}$--$150^{\circ}$ (DoE) \\
            & Injection side       & \textbf{Bottom} ($y = -R_1$) \\
Symmetry    & Plane                & $x = 0$ (half-domain, $x \geq 0$) \\
\bottomrule
\end{tabular}
\end{table}

\subsection{Bottom-Injection Parameterisation}

The \texttt{generateSTL.py} script uses two environment variables to
control the injection geometry:

\begin{itemize}
  \item \texttt{BOTTOM\_INJECT=1} flips the branch to $y = -R_1$.
        Internally, a sign variable \texttt{Y\_SIGN = -1} is applied
        to the junction vector, branch direction vector $\hat{s}$,
        hole-detection logic, and intersection curve calculation.
  \item \texttt{ALPHA\_DEG} sets the injection angle (degrees from
        the $-z$ main-flow axis).  $90^{\circ}$ is perpendicular;
        $30^{\circ}$ is nearly co-flow; $150^{\circ}$ is
        counter-flow.
\end{itemize}


\section{Step 1: Generate the STL}

\begin{lstlisting}[language=bash,numbers=none]
source case.env          # sets D2, ALPHA_DEG, BOTTOM_INJECT
python3 generateSTL.py   # writes constant/triSurface/*.stl
\end{lstlisting}

Four named STL surfaces are created:
\begin{enumerate}
  \item \texttt{wall.stl} --- combined pipe walls (watertight).
  \item \texttt{main\_inlet.stl} --- cap at $Z = 0$.
  \item \texttt{outlet.stl} --- cap at $Z = L$.
  \item \texttt{branch\_inlet.stl} --- cap at the branch tip.
\end{enumerate}

The script performs a signed-volume sanity check (target $\pm 10\%$
of the analytic volume) and aborts if the geometry is not watertight.


\section{Step 2: Background Mesh (\texttt{blockMesh})}

The half-domain background mesh extends from $x = 0$ (symmetry plane)
to $x = 0.30$\,m:

\begin{table}[H]
\centering
\caption{\texttt{blockMeshDict} extents (angle-dependent $y_{\min}$).}
\begin{tabular}{@{}lrrr@{}}
\toprule
\textbf{Direction} & \textbf{Min (m)} & \textbf{Max (m)} & \textbf{Cells} \\
\midrule
$X$ & $0.00$ & $+0.30$ & 10 \\
$Y$ ($90^{\circ}$) & $-3.10$ & $+0.30$ & 115 \\
$Y$ ($30^{\circ}$) & $-1.80$ & $+0.30$ & 70 \\
$Z$ & $-0.15$ & $+6.95$ & 237 \\
\bottomrule
\end{tabular}
\end{table}

The $y_{\min}$ is computed from the worst-case branch tip position
$y_{\mathrm{tip}} = -(R_1 + L_{\mathrm{branch}} \sin\alpha)$ with a
0.15\,m buffer.


\section{Steps 3--4: Features + \texttt{snappyHexMesh}}

\begin{lstlisting}[language=bash,numbers=none]
blockMesh
surfaceFeatureExtract
snappyHexMesh -overwrite
\end{lstlisting}

\noindent The \texttt{snappyHexMeshDict} performs:

\begin{itemize}
  \item \textbf{Castellate:} refinement level 2 on pipe walls,
        level 3 inside a junction body-of-influence sphere
        (radius $\sim 3\,D_2$).
  \item \textbf{Snap:} implicit + explicit feature snap (10 iterations).
  \item \textbf{Layers:} 3 prismatic layers on walls, expansion ratio
        1.2.
\end{itemize}

Target cell count: \textbf{400--500\,K} (half-domain RANS).

\subsection{Boundary Patches}

\begin{table}[H]
\centering
\caption{Patches created by \texttt{snappyHexMesh}.}
\begin{tabular}{@{}lll@{}}
\toprule
\textbf{Patch} & \textbf{Type} & \textbf{Description} \\
\midrule
\texttt{main\_inlet}   & \texttt{patch} & $Z = 0$ cap \\
\texttt{branch\_inlet}  & \texttt{patch} & Branch tip (at $y < 0$) \\
\texttt{outlet}         & \texttt{patch} & $Z = L$ cap \\
\texttt{wall}           & \texttt{wall}  & All pipe walls \\
\texttt{sym}            & \texttt{symmetryPlane} & $x = 0$ plane \\
\bottomrule
\end{tabular}
\end{table}

\subsection{Known Meshing Limitation}

For the smallest branch diameter ($d/D \approx 0.196$),
\texttt{snappyHexMesh} occasionally produces a \textbf{disconnected
mesh} (2 regions), preventing H$_2$ from entering the main pipe.
This occurs because the hole in the main pipe wall is too small
for the castellated background mesh to resolve.  Cases with $d/D
\geq 0.25$ produce single-region meshes reliably.


% ════════════════════════════════════════════════════════════════
%  CHAPTER 4 — Physical Models and Properties
% ════════════════════════════════════════════════════════════════
\chapter{Physical Models and Properties}
\label{ch:models}

\section{Turbulence Model}

\begin{lstlisting}[language={},title=constant/momentumTransport]
simulationType  RAS;

RAS
{
    RASModel        kOmegaSST;
    turbulence      on;
    printCoeffs     on;
}
\end{lstlisting}

\begin{lstlisting}[language={},title=constant/turbulenceProperties]
simulationType  RAS;
RAS
{
    RASModel        kOmegaSST;
    turbulence      on;
    printCoeffs     on;
}
Sct     0.5;
kMin    1e-10;
omegaMin 1e-10;
\end{lstlisting}

\noindent The $k$-$\omega$ SST model was chosen over DES/DDES for
computational efficiency in the parametric study (10--15 cases per
angle).  With the shortened domain and half-symmetry, each case
completes in 40--80 minutes on 20 cores.

\section{Thermophysical Properties}

The binary H$_2$/CH$_4$ mixture uses JANAF thermodynamics,
Sutherland transport, and the perfect-gas equation of state:

\begin{lstlisting}[language={},title=constant/thermophysicalProperties (header)]
thermoType
{
    type            heRhoThermo;
    mixture         reactingMixture;
    transport       sutherland;
    thermo          janaf;
    energy          sensibleEnthalpy;
    equationOfState perfectGas;
    specie          specie;
}
inertSpecie CH4;
species     ( H2 CH4 );
\end{lstlisting}

\begin{table}[H]
\centering
\caption{Species properties at 6.9\,MPa, 288\,K.}
\label{tab:species}
\small
\begin{tabular}{@{}lrrl@{}}
\toprule
\textbf{Property} & \textbf{CH$_4$} & \textbf{H$_2$} & \textbf{Source} \\
\midrule
Molecular weight (kg/kmol) & 16.043  & 2.016 & NIST \\
Density at 6.9\,MPa (kg/m$^3$) & $\sim 45$ & $\sim 5.5$ & ideal gas \\
Dynamic viscosity (Pa\,s) & $1.25\!\times\!10^{-5}$ & $9.0\!\times\!10^{-6}$ & NIST \\
Sutherland $A_s$ (Pa\,s/K$^{1/2}$) & $1.17\!\times\!10^{-6}$ & $6.69\!\times\!10^{-7}$ & fitted \\
Sutherland $T_s$ (K) & 164 & 72 & standard \\
\bottomrule
\end{tabular}
\end{table}


\section{Chemistry \& Combustion (Disabled)}

\begin{lstlisting}[language={},title=constant/combustionProperties]
combustionModel none;
\end{lstlisting}

\begin{lstlisting}[language={},title=constant/reactions]
species   ( H2 CH4 );
reactions ();
\end{lstlisting}


\section{Gravity}

\begin{lstlisting}[language={},title=constant/g]
dimensions  [0 1 -2 0 0 0 0];
value       (0 -9.81 0);
\end{lstlisting}

\noindent Gravity acts along $-Y$.  For bottom injection ($y = -R_1$),
the lighter H$_2$ experiences a net upward buoyancy force, causing the
plume centroid to \textbf{rise} as it convects downstream.  This was
verified quantitatively: the H$_2$-weighted $y$-centroid increases
monotonically from $y \approx -0.06$\,m at $z = 3$\,m to
$y \approx +0.03$\,m at $z = 6$\,m for the highest velocity-ratio case
(VR\,=\,3.81).


% ════════════════════════════════════════════════════════════════
%  CHAPTER 5 — Boundary and Initial Conditions
% ════════════════════════════════════════════════════════════════
\chapter{Boundary and Initial Conditions}
\label{ch:bc}

\section{Power-Law Inlet Velocity Profile}

Both inlets use a $1/7$-power-law turbulent pipe profile implemented
via \texttt{codedFixedValue}:

\begin{equation}
  u(r) = u_{\max}\!\left[1 - \left(\frac{2r}{d}\right)^{\!2}\right]^{1/7}
  \label{eq:powerlaw}
\end{equation}

\subsection{Main Inlet}

The main inlet velocity is $U_{\mathrm{main}} = 10$\,m/s along $+Z$,
with the radial coordinate measured from the pipe centreline at
$(x, y) = (0, 0)$.

\subsection{Branch Inlet (Bottom Injection)}

The branch inlet uses the same profile but oriented along the
\emph{negative} branch-axis direction $-\hat{s}$, where:
\[
  \hat{s} = (0,\; \texttt{Y\_SIGN} \cdot \sin\alpha,\; -\cos\alpha)
\]

For bottom injection (\texttt{Y\_SIGN}\,=\,$-1$, $\alpha = 90^{\circ}$),
the velocity vector at each face is:
\[
  \mathbf{u} = u(r) \cdot (0,\; +\sin\alpha,\; \cos\alpha)
\]
i.e.\ the jet enters upward ($+y$) into the main pipe.  The
\texttt{@Y\_SIGN@} token is substituted by \texttt{stamp\_cases.py}
at case generation time.


\section{Pressure}

\begin{table}[H]
\centering
\caption{Pressure boundary conditions.}
\begin{tabularx}{\textwidth}{@{}lllX@{}}
\toprule
\textbf{Patch} & \textbf{\texttt{p\_rgh} type} & \textbf{Value} & \textbf{Notes} \\
\midrule
\texttt{main\_inlet}   & \texttt{fixedFluxPressure} & --- & Balances body force \\
\texttt{branch\_inlet}  & \texttt{fixedFluxPressure} & --- & Balances body force \\
\texttt{outlet}         & \texttt{prghPressure} & $6.9 \times 10^6$\,Pa & Hydrostatic removed \\
\texttt{wall}           & \texttt{fixedFluxPressure} & --- & Balances body force \\
\texttt{sym}            & \texttt{symmetryPlane} & --- & \\
\bottomrule
\end{tabularx}
\end{table}


\section{Temperature, Species, Turbulence}

\begin{table}[H]
\centering
\caption{Scalar boundary conditions.}
\small
\begin{tabularx}{\textwidth}{@{}lllllll@{}}
\toprule
\textbf{Patch} & $T$ & $Y_{\mathrm{H_2}}$ & $k$ & $\omega$ & $\nu_t$ \\
\midrule
\texttt{main\_inlet}   & 288\,K & 0 & calc & calc & calc \\
\texttt{branch\_inlet}  & 288\,K & 1 & calc & calc & calc \\
\texttt{outlet}         & \texttt{inletOutlet} & \texttt{inletOutlet} & \texttt{inletOutlet} & \texttt{inletOutlet} & calc \\
\texttt{wall}           & \texttt{zeroGradient} & \texttt{zeroGradient} & \texttt{kqRWallFunction} & \texttt{omegaWallFunction} & \texttt{nutkWallFunction} \\
\texttt{sym}            & \texttt{symmetryPlane} & \texttt{symmetryPlane} & \texttt{symmetryPlane} & \texttt{symmetryPlane} & \texttt{symmetryPlane} \\
\bottomrule
\end{tabularx}
\end{table}

\noindent An \texttt{allBoundary} fallback entry is included in all
boundary-condition files to prevent crashes if \texttt{snappyHexMesh}
leaves a residual background-mesh patch.


% ════════════════════════════════════════════════════════════════
%  CHAPTER 6 — Solution Methods and Controls
% ════════════════════════════════════════════════════════════════
\chapter{Solution Methods and Controls}
\label{ch:solution}

\section{Pressure--Velocity Coupling}

\begin{lstlisting}[language={},title=system/fvSolution (PIMPLE block)]
PIMPLE
{
    momentumPredictor        yes;
    nOuterCorrectors         1;      // = PISO
    nCorrectors              2;
    nNonOrthogonalCorrectors 1;
    transonic                no;
}
\end{lstlisting}


\section{Discretisation Schemes}

\begin{table}[H]
\centering
\caption{Convection schemes (first-order upwind for RANS robustness).}
\label{tab:divschemes}
\begin{tabularx}{\textwidth}{@{}lX@{}}
\toprule
\textbf{Term} & \textbf{Scheme} \\
\midrule
\texttt{div(phi,U)}       & \texttt{Gauss upwind} \\
\texttt{div(phi,H2)}      & \texttt{Gauss upwind} \\
\texttt{div(phi,k)}       & \texttt{Gauss upwind} \\
\texttt{div(phi,omega)}   & \texttt{Gauss upwind} \\
\texttt{div(phi,h)}       & \texttt{Gauss upwind} \\
Time (\texttt{ddtSchemes}) & \texttt{Euler} (with adaptive $\Delta t$) \\
Gradient                   & \texttt{Gauss linear} \\
Laplacian                  & \texttt{Gauss linear corrected} \\
\bottomrule
\end{tabularx}
\end{table}

\noindent First-order upwind introduces $\sim 5$--$10\%$ spatial
diffusion error on CoV (Ferziger \& Peri\'{c} 2002), which is
acceptable for the comparative parametric study where relative trends
between cases are the primary output.


\section{Time-Stepping}

\begin{table}[H]
\centering
\caption{Run parameters.}
\begin{tabular}{@{}lr@{}}
\toprule
\textbf{Parameter} & \textbf{Value} \\
\midrule
Application            & \texttt{rhoReactingBuoyantFoam} \\
Initial $\Delta t$     & $10^{-5}$\,s \\
Adaptive time-stepping & Yes (\texttt{adjustTimeStep}) \\
Max Courant number     & 4.5 \\
End time               & 1.2\,s ($\approx 2$ flow-throughs) \\
Write interval         & 0.1\,s \\
Purge write            & 3 (keep last 3 snapshots) \\
Parallel decomposition & 20 subdomains (scotch) \\
\bottomrule
\end{tabular}
\end{table}


% ════════════════════════════════════════════════════════════════
%  CHAPTER 7 — Function Objects and Monitors
% ════════════════════════════════════════════════════════════════
\chapter{Function Objects and Monitors}
\label{ch:monitors}

All monitors are declared in \texttt{controlDict} under
\texttt{functions}; output lands in \texttt{postProcessing/}.

\section{Time-Averaged Fields}

\begin{lstlisting}[language={},numbers=none]
fieldAverage1
{
    type            fieldAverage;
    writeControl    writeTime;
    timeStart       0.6;  // start after 1 flow-through
    fields
    (
        U    { mean on; prime2Mean on; }
        H2   { mean on; prime2Mean on; }
        p_rgh { mean on; prime2Mean off; }
        k    { mean on; prime2Mean off; }
    );
}
\end{lstlisting}

This produces \texttt{UMean}, \texttt{H2Mean}, \texttt{p\_rghMean},
etc.\ in the final time directory.

\section{Sampling Planes}

Area-averaged H$_2$ and $p_{\mathrm{rgh}}$ are sampled at four
downstream cross-sections:

\begin{table}[H]
\centering
\caption{Sampling plane locations.}
\begin{tabular}{@{}lrr@{}}
\toprule
\textbf{Plane} & $z$ (m) & Distance from junction \\
\midrule
\texttt{planeZ3} & 3.0 & $0.7$\,m ($\approx 1.5\,D$) \\
\texttt{planeZ4} & 4.0 & $1.7$\,m ($\approx 3.7\,D$) \\
\texttt{planeZ5} & 5.0 & $2.7$\,m ($\approx 5.9\,D$) \\
\texttt{planeZ6} & 6.0 & $3.7$\,m ($\approx 8.0\,D$) \\
\texttt{outlet}  & 6.9 & $4.6$\,m ($= 10\,D$) \\
\bottomrule
\end{tabular}
\end{table}

\section{Mass Balance}

\texttt{surfaceFieldValue} function objects integrate mass flux
$\phi$ over \texttt{main\_inlet} and \texttt{outlet}.  The
time-averaged outlet mass flow should match the sum of inlets to
within $< 5\%$.

\section{Residuals and $y^+$}

\begin{itemize}
  \item \texttt{residuals} --- per-timestep residuals for
        $U$, $p_{\mathrm{rgh}}$, $H_2$, $k$, $\omega$.
  \item \texttt{yPlusAudit} --- wall $y^+$ at write times.
\end{itemize}


% ════════════════════════════════════════════════════════════════
%  CHAPTER 8 — Design of Experiments
% ════════════════════════════════════════════════════════════════
\chapter{Design of Experiments}
\label{ch:doe}

\section{DoE Variables}

\begin{table}[H]
\centering
\caption{DoE parameter ranges.}
\begin{tabular}{@{}llrr@{}}
\toprule
\textbf{Variable} & \textbf{Symbol} & \textbf{Min} & \textbf{Max} \\
\midrule
Branch diameter ratio & $d/D$ & 0.196 & 0.45 \\
Hydrogen blend ratio  & HBR   & 0.06  & 0.20 \\
Injection angle       & $\alpha$ & $30^{\circ}$ & $150^{\circ}$ \\
\bottomrule
\end{tabular}
\end{table}

\noindent The velocity ratio VR is derived from $(d/D, \mathrm{HBR})$
and the density ratio.  Each angle is run as a separate campaign of
10 cases using Sliced Latin Hypercube Sampling (5 $d/D$ slices
$\times$ 2 HBR points per slice).

\section{Sliced LHS Design}

The slicing ensures that each pair of cases within a $d/D$ slice
shares the same \texttt{snappyHexMesh} result (mesh reuse), reducing
the total meshing cost by $\sim 50\%$.  Cases are sorted by expected
wall time (lowest Re first) for efficient sequential execution.

\section{Case Stamping (\texttt{stamp\_cases.py})}

The stamper reads \texttt{doe\_design.csv} and for each row:
\begin{enumerate}
  \item Copies the base template to \texttt{case\_NN/}.
  \item Replaces tokens (\texttt{@D2@}, \texttt{@UBY@},
        \texttt{@Y\_SIGN@}, etc.) in all \texttt{0/} and
        \texttt{system/} files.
  \item Writes \texttt{case.env} (shell-sourceable parameters) and
        \texttt{case\_info.json} (machine-readable metadata).
\end{enumerate}


% ════════════════════════════════════════════════════════════════
%  CHAPTER 9 — Results: 90° Bottom Injection
% ════════════════════════════════════════════════════════════════
\chapter{Results: $90^{\circ}$ Bottom Injection}
\label{ch:results90}

\section{Campaign Summary}

\begin{table}[H]
\centering
\caption{90° bottom-injection results (8 valid cases; cases 01--02
         excluded due to disconnected mesh at $d/D = 0.196$).
         CoV and $\Delta p$ from function-object time series.}
\label{tab:results90}
\small
\begin{tabular}{@{}rrrrrrr@{}}
\toprule
\textbf{Case} & $d/D$ & VR & $\langle Y_{\mathrm{H_2}}\rangle_{\mathrm{outlet}}$ (\%) & \textbf{CoV} (area) & $\Delta p$ (kPa) & Wall time (s) \\
\midrule
03 & 0.252 & 1.33 & 1.04 & 1.10 & 3.91 & 2387 \\
04 & 0.252 & 3.81 & 3.26 & \textbf{0.27} & 0.27 & 4900 \\
05 & 0.296 & 1.24 & 1.28 & 0.99 & 4.68 & 2764 \\
06 & 0.296 & 2.61 & 2.19 & 0.95 & 3.95 & 4448 \\
07 & 0.382 & 1.14 & 1.83 & 0.95 & 3.68 & 2926 \\
08 & 0.382 & 0.69 & 0.23 & 0.91 & $-1.31$ & 2273 \\
09 & 0.396 & 0.95 & 1.09 & 0.83 & 0.73 & 2745 \\
10 & 0.396 & 0.81 & 0.52 & 0.84 & $-0.92$ & 2390 \\
\bottomrule
\end{tabular}
\end{table}

\subsection{Supplementary High-VR Cases}

Five additional cases with high velocity ratios were run to explore
the regime where the H$_2$ plume crosses the pipe centreline:

\begin{table}[H]
\centering
\caption{Supplementary cases targeting strong plume rise.}
\small
\begin{tabular}{@{}rrrrl@{}}
\toprule
\textbf{Case} & $d/D$ & VR & $U_{\mathrm{branch}}$ (m/s) & \textbf{Rationale} \\
\midrule
11 & 0.350 & 4.0 & 40 & Medium branch + very high VR \\
12 & 0.400 & 3.5 & 35 & Large branch + high VR \\
13 & 0.450 & 3.0 & 30 & Max branch + high VR \\
14 & 0.450 & 4.5 & 45 & Max branch + very high VR \\
15 & 0.350 & 5.0 & 50 & Extreme VR, deepest penetration \\
\bottomrule
\end{tabular}
\end{table}

\section{Key Physical Findings}

\subsection{Mixing Effectiveness}

\begin{itemize}
  \item \textbf{Higher VR $\Rightarrow$ lower CoV:}  Across all $d/D$
        pairs, the higher-VR case consistently achieves better outlet
        mixing.  Case~04 (VR\,=\,3.81) achieves CoV\,=\,0.27
        (best in campaign).
  \item \textbf{Larger $d/D$ improves mixing:}  At comparable VR,
        larger branch diameters produce lower CoV due to greater
        contact surface area.
\end{itemize}

\subsection{Buoyancy Verification}

The H$_2$-weighted $y$-centroid at downstream cross-sections confirms
that buoyancy is acting correctly:

\begin{table}[H]
\centering
\caption{H$_2$ plume centroid $y$-position (m) at downstream stations.}
\small
\begin{tabular}{@{}lrrrrl@{}}
\toprule
\textbf{Case} & $z = 3$\,m & $z = 4$\,m & $z = 5$\,m & $z = 6$\,m & \textbf{Rising?} \\
\midrule
03 (VR = 1.33) & $-0.175$ & $-0.151$ & $-0.132$ & $-0.148$ & Yes \\
04 (VR = 3.81) & $-0.056$ & $-0.001$ & $+0.019$ & $+0.031$ & \textbf{Yes} \\
07 (VR = 1.14) & $-0.164$ & $-0.134$ & $-0.120$ & $-0.137$ & Yes \\
09 (VR = 0.95) & $-0.161$ & $-0.137$ & $-0.121$ & $-0.123$ & Yes \\
\bottomrule
\end{tabular}
\end{table}

\noindent All cases show the plume rising from the bottom toward
the pipe centre, with case~04 crossing the centreline by
$z = 4$\,m.  This confirms:
\begin{enumerate}
  \item Gravity is correctly oriented ($g_y = -9.81$\,m/s$^2$).
  \item The perfect-gas density coupling produces the correct
        buoyancy force on the lighter H$_2$.
  \item Higher VR produces deeper jet penetration \emph{and}
        stronger buoyancy-assisted rise.
\end{enumerate}

\subsection{Pressure Drop}

The time-series-averaged $\Delta p$ ranges from $-1.3$ to $+4.7$\,kPa.
Low-VR cases (08, 10) show slightly negative $\Delta p$, indicating
the branch barely perturbs the main flow.  Higher VR cases produce
positive $\Delta p$ from the momentum disruption at the junction.


\section{Figures Generated}

For each case, the following figures are produced automatically by
\texttt{make\_figures.py} and \texttt{make\_distance\_figures.py}:

\begin{table}[H]
\centering
\caption{Figure inventory per case.}
\small
\begin{tabularx}{\textwidth}{@{}lX@{}}
\toprule
\textbf{Figure} & \textbf{Description} \\
\midrule
\texttt{fig\_H2\_long.png}    & Longitudinal centreline slice of H$_2$Mean \\
\texttt{fig\_H2\_strip.png}   & Cross-section strip at $z = 3, 4, 5, 6$\,m + outlet \\
\texttt{fig\_H2\_isosurface.png} & 3D isosurface of H$_2$ plume \\
\texttt{fig\_H2\_cross\_sections.png} & 3D view of cross-sections in the pipe \\
\texttt{fig\_H2\_topdown.png}  & Top-down view of H$_2$ distribution \\
\texttt{fig\_H2\_outlet.png}   & H$_2$ concentration at the outlet face \\
\texttt{fig\_velocity\_xz.png} & Velocity magnitude on centreline slice \\
\texttt{fig\_pressure\_xz.png} & $p_{\mathrm{rgh}}$ on centreline slice \\
\texttt{fig\_mesh\_xz.png}    & Mesh cells on centreline slice \\
\texttt{fig\_geometry.png}     & STL geometry (isometric view) \\
\bottomrule
\end{tabularx}
\end{table}


% ════════════════════════════════════════════════════════════════
%  CHAPTER 10 — Post-Processing
% ════════════════════════════════════════════════════════════════
\chapter{Post-Processing}
\label{ch:post}

\section{Metrics Extraction (\texttt{all\_metrics.py})}

The retro-active post-processor computes from reconstructed boundary
fields:

\begin{itemize}
  \item \textbf{Mixing at outlet:} area-, mass-flux-, and
        volume-flux-weighted $\langle Y_{\mathrm{H_2}} \rangle$,
        $\sigma(Y_{\mathrm{H_2}})$, CoV, and Danckwerts intensity
        of segregation $I_s$.
  \item \textbf{Pressure drop:} $\Delta p$ for static, gauge,
        total, and gauge-total pressures, each with three weightings
        (12 values).
  \item \textbf{Mass balance:} direct integration of $\dot{m}$ at
        all three patches, with closure error percentage.
\end{itemize}

\begin{lstlisting}[language=bash,numbers=none]
python3 tools/all_metrics.py case_03 --outdir results/case_03
\end{lstlisting}

Outputs: \texttt{all\_metrics.csv} and \texttt{ALL\_METRICS.md}.

\section{Coefficient of Variation}

\[
  \mathrm{CoV} = \frac{\sigma_{Y_{\mathrm{H_2}}}}
                       {\langle Y_{\mathrm{H_2}} \rangle}
\]

$\mathrm{CoV} = 0$ denotes perfect mixing; $\mathrm{CoV} > 1$
denotes severe segregation.  The preferred computation uses the
\textbf{time-averaged field} \texttt{H2Mean} to remove acoustic
noise from the compressible solver.

\section{Figure Generation}

\begin{lstlisting}[language=bash,numbers=none]
pip install pyvista numpy matplotlib
PYVISTA_OFF_SCREEN=true python3 tools/make_figures.py case_03 case_03/figures
PYVISTA_OFF_SCREEN=true python3 tools/make_distance_figures.py \
    --case case_03 --outdir case_03/figures
\end{lstlisting}

The scripts use PyVista for 3D rendering with a triangulated
volume mesh to avoid white-notch artefacts from polyhedral
non-manifold cells at the junction.


% ════════════════════════════════════════════════════════════════
%  CHAPTER 11 — Workflow and Campaign Management
% ════════════════════════════════════════════════════════════════
\chapter{Workflow and Campaign Management}
\label{ch:workflow}

\section{End-to-End Pipeline}

\begin{enumerate}
  \item \textbf{Generate DoE design:}
\begin{lstlisting}[language=bash,numbers=none]
ALPHA_DEG=90 python3 doe/lhs_design.py --alpha-deg 90 --out doe/doe_design.csv
\end{lstlisting}

  \item \textbf{Stamp cases:}
\begin{lstlisting}[language=bash,numbers=none]
ALPHA_DEG=90 BOTTOM_INJECT=1 python3 doe/stamp_cases.py \
    --design doe/doe_design.csv --base . --cases doe_cases
\end{lstlisting}

  \item \textbf{Clean \texttt{cellLevel} from all cases:}
\begin{lstlisting}[language=bash,numbers=none]
for c in doe_cases/case_*/; do
    rm -f "$c/0/cellLevel" "$c/0/pointLevel"
    rm -rf "$c/constant/polyMesh"
done
\end{lstlisting}

  \item \textbf{Launch campaign:}
\begin{lstlisting}[language=bash,numbers=none]
tmux new-session -d -s doe
tmux send-keys -t doe \
    'CASES_DIR=doe_cases RESULTS_DIR=doe_results \
     TOOLS_DIR=tools bash doe/run_doe.sh 2>&1 | tee run.log' Enter
\end{lstlisting}
\end{enumerate}

\section{Known Pitfalls}

\begin{enumerate}
  \item \textbf{Stale \texttt{cellLevel}:}  When
        \texttt{stamp\_cases.py} copies the base, it may include a
        residual \texttt{0/cellLevel} from a previous
        \texttt{snappyHexMesh} run.  If mesh-reuse then copies a
        \texttt{polyMesh} with a different cell count,
        \texttt{decomposePar} crashes with ``Size $N$ is not equal
        to the expected length $M$''.  \textbf{Always} clean
        \texttt{cellLevel} and \texttt{pointLevel} from all cases
        before launching.
  \item \textbf{Stale \texttt{polyMesh} in base:}  Ensure
        \texttt{constant/polyMesh} does not exist in the base
        template before stamping, or it will be copied into every
        case and interfere with the mesh-donor logic.
  \item \textbf{Disconnected mesh at small $d/D$:}  For $d/D
        \lesssim 0.20$, the branch opening is too small for the
        background mesh to resolve.  Mark these cases as invalid
        rather than re-deriving the DoE.
  \item \textbf{No internet on compute nodes:}  Install Python
        dependencies (\texttt{pyvista}, \texttt{matplotlib}) on a
        machine with internet access, then \texttt{rsync} the data
        for local post-processing.
\end{enumerate}


\vfill
\noindent\fbox{\parbox{\dimexpr\textwidth-2\fboxsep-2\fboxrule}{%
\textbf{Pre-flight checklist:}
\begin{itemize}[nosep]
  \item[\checkmark] OpenFOAM 2406 environment sourced.
  \item[\checkmark] \texttt{case.env} contains correct \texttt{ALPHA\_DEG},
        \texttt{BOTTOM\_INJECT=1}, \texttt{D2}.
  \item[\checkmark] \texttt{0/cellLevel} and \texttt{0/pointLevel}
        deleted from all cases.
  \item[\checkmark] No stale \texttt{constant/polyMesh} in base template.
  \item[\checkmark] \texttt{decomposeParDict} \texttt{numberOfSubdomains}
        matches \texttt{mpirun -np} count.
  \item[\checkmark] \texttt{blockMeshDict} $y_{\min}$ covers worst-case
        branch extent for the chosen angle.
  \item[\checkmark] \texttt{controlDict} \texttt{application}
        = \texttt{rhoReactingBuoyantFoam}.
\end{itemize}}}

\end{document}
